%% kapa: exam=1 round=23 year=2012 period=1 form=A subject=civil_law subject_ko=민법 seq=25 number=25
%% source: Q-Net 감정평가사 기출문제 (cst00309) artlSeq=5052578 file=raw/1008/5052578/2012년 제23회 감정평가사 제1차시험 기출문제(1교시 A형).pdf
%% conversion: kordoc 4.24.51 -> tools/gen1.py (auto); review=proofread-sonnet; flags=none
\begin{question}{25}
甲이 등기를 하여야 비로소 소유권을 취득하는 경우는? (다툼이 있으면 판례에 의함)
\begin{choices}
  \item 甲이 수용재결을 통해 토지를 취득한 경우
  \item 甲 소유의 건물에 임차인이 창고를 지어 부합된 경우
  \item 저당권 실행을 위한 경매에서 매수인 甲이 매각대금을 완납한 경우
  \item 乙이 건물을 신축한 후 사망하여 甲이 상속한 경우
  \item 甲이 부동산을 20년간 평온ㆍ공연하게 소유의 의사로 점유하여 점유취득시효를 완성한 경우
\end{choices}
\end{question}
\answer{5}
